\section{The matched cubic cost after retaining one weight}
\label{sec:mixed-cubic-residual-trivial-window}

\paragraph{The unchanged datum and retained arithmetic sum.}
Retain every original field, coefficient, lattice, full local condition,
normalization, finite datum, page and prescribed pair of AC-G218-window.
In particular $h\geq2$ and $T\geq2$ are integers, $N=T-1$, and
\[
 \begin{gathered}
 H_u=\mathbb Q(\zeta_{2^{h+2}},\zeta_7,\sqrt{-q},\sqrt{-d},
                         \omega,\sqrt[3]{269}),\qquad K=\mathbb Q(\omega),\\
 u^{2+8\varepsilon}\leq2^h\leq2u^{2+8\varepsilon},\qquad
 T=\lfloor u^{399}\rfloor,\qquad 0<\varepsilon<1/256.
 \end{gathered}\tag{MRTW.1}
\]
Keep the complete old sets $Z_u,V_u,Y_u,C_5^{\rm all},C_7^{\rm all}$,
$U_5^{\rm all},U_7^{\rm all}$ and $M_{C7}$ of (MQSW.3)--(MQSW.8).
Abbreviate the whole finite nonnegative remainder (MQSW.5) by
$\mathcal R^{\rm old}(u)$, and put
\[
 \begin{split}
 J_u^*&=(J_{u,\rm mix}\setminus M_{C7})\ \sqcup\
       \bigl(J_{u,0}\setminus(V_u\cup Y_u\cup C_5^{\rm all}
                                             \cup C_7^{\rm all})\bigr),\\
 \mathcal R^{\rm old}(u)&=\mathcal R_{J_u^*}.
 \end{split}\tag{MRTW.2}
\]
The right side uses exactly the mixed and noncritical summands of
(CP.3), on their original rings. No previously accounted-for triple
is restored. Define the full residual-matched mixed family
\[
 \begin{split}
 \rho_{a,e}&=\bar\xi_7^{\,a}\bar c_7^{-e},\\
 B_u=\{(7,\chi^j\xi_7^a,e):{}&j=1,2,\quad0\leq a\leq5,
               \quad1\leq e<T,\quad\rho_{a,e}=1\}.
 \end{split}\tag{MRTW.3}
\]
The rational two-power and auxiliary $q,d$ factors are trivial.
The residual equality is literal in the prescribed original embedding.
For each $m\in\{1,\ldots,N\}$ let
\[
 \begin{split}
 P_m(u)&=\{(7,\eta,e)\in B_u:e=m\},\qquad
 A_m(u)=B_u\setminus P_m(u),\\
 \mathcal R_{399}^{[m]}(u)&=\mathcal R_{J_u^*\setminus A_m(u)},\\
 \mathcal R_{399}^{\rm matched,max}(u)&=
                   \max_{1\leq m\leq N}\mathcal R_{399}^{[m]}(u),\\
 \mathcal R^{\rm core}(u)&=\mathcal R_{J_u^*\setminus B_u}.
 \end{split}\tag{MRTW.4}
\]
Every branch retains the pair of matched cubic labels at one
weight, the entire old rational noncritical sum, and every mixed
term outside $B_u$. The core is used only to express this identity;
no sufficient numerical gate is claimed for the core by itself.
The maximum is a completely specified function of the same actual
datum and the finitely many weights $1,\ldots,N$.

\begin{theorem}[AC-G220-window: the matched cubic cost with one retained weight]
\label{AC-G220-window}
The higher-lifting conclusion (MRT.27)--(MRT.31) of AC-G219
supplies integers $s_e\geq1$, independent of the original coefficient
ring, host and two cubic labels, such that every matched group is
$O_7/7^{s_e}O_7$. For each $n\geq2$, the set of positive weights
with $s_e\geq n$ is contained in at most one residue class modulo
$7^{n-1}$, with compatible classes when they exist.
The exhaustive old disjointness and counts are
\[
 \begin{gathered}
 B_u\subset J_{u,\rm mix}\setminus M_{C7},\qquad |B_u|=2N,\\
 B_u\cap(Z_u\cup V_u\cup Y_u\cup C_5^{\rm all}
                    \cup C_7^{\rm all}\cup M_{C7})=\varnothing,\\
 |P_m(u)|=2,\qquad |A_m(u)|=2(N-1).
 \end{gathered}\tag{MRTW.5}
\]
In particular,
\[
 \begin{aligned}
 \mathcal R^{\rm old}&=\mathcal R^{\rm core}
                             +2\log7\sum_{e=1}^{N}s_e,\\
 \mathcal R_{399}^{[m]}&=\mathcal R^{\rm core}+2s_m\log7,\\
 \mathcal R_{399}^{\rm matched,max}
       &=\mathcal R^{\rm core}+2\log7\max_{1\leq e\leq N}s_e.
 \end{aligned}\tag{MRTW.6}
\]
Define the explicit finite integer
\[
 L_N=(N-1)+\sum_{k\geq1}\left\lfloor\frac{N-1}{7^k}\right\rfloor
       =(N-1)+v_7((N-1)!).
 \tag{MRTW.7}
\]
Then the exact nonnegative comparison cost obeys
\[
 \begin{gathered}
 \Gamma_u:=\mathcal R^{\rm old}
                       -\mathcal R_{399}^{\rm matched,max}\\
       =2\log7\left(\sum_{e=1}^{N}s_e-\max_{1\leq e\leq N}s_e\right),\\
 2(N-1)\log7\leq\Gamma_u\leq2L_N\log7\\
       \leq\frac73(N-1)\log7<7\mathcal P_u(T).
 \end{gathered}\tag{MRTW.8}
\]
Every branch indexing set is properly reduced for $T\geq3$,
and the maximum is strictly smaller than the old finite remainder
there. At $T=2$ it equals the old remainder.
No fixed maximizing weight, or upper bound for that retained
weight's group, is assumed.

For the old growing-prime count put
\[
 \mathcal N_u^{\rm grow}=8J_{\rm cub}(T)+4J_7(T)+4J_C(T)
                       +2\sum_{D\in\{-7,7,-14,14\}}J_D(T).
 \tag{MRTW.9}
\]
For every branch the total represented-triple count outside its
remainder is
\[
 \begin{gathered}
 |Z_u|+|V_u|+|Y_u|+67N-2-\lfloor N/4\rfloor
                                      +\mathcal N_u^{\rm grow}.
 \end{gathered}\tag{MRTW.10}
\]
Thus the count is independent of which weight is retained.
The maximum is not identified with one fixed deletion set, and
the $N$ branch counts are not added together.

Retain exactly the cost functions $\mathcal E_u,\mathcal L_u,
\mathcal A_u,\mathcal P_u$ and absolute constant $C_2$ of
(CAW.5) and (MQSW.9). One has
\[
 \begin{aligned}
 |I_{399}-\mathcal R_{399}^{\rm matched,max}|&\leq
     \mathcal E_u+\mathcal L_u+\mathcal A_u
                                      +(C_2+7)\mathcal P_u,\\
 |I_{399}^{\rm all7,2}-\mathcal R_{399}^{\rm matched,max}|&\leq
     \mathcal E_u+\mathcal L_u+(C_2+7)\mathcal P_u.
 \end{aligned}\tag{MRTW.11}
\]
The sharper additional cost is $\Gamma_u$, or its explicit upper
bound $2L_N\log7$, instead of $7\mathcal P_u$.
There is still one full finite-gluing error, one degree-one torsion
aggregate, one old critical aggregate and the preceding G189 cost
when the full index is used. The added matched cost is charged once.
The uniform error envelope remains
\[
 O_\varepsilon\bigl(u^{806+16\varepsilon}(\log u)^2\bigr)
                         =o\bigl(\exp(u^\varepsilon/2)\bigr).
 \tag{MRTW.12}
\]
The sufficient arithmetic premise uses the identical prescribed
window with the adjusted threshold and slack specified below.
The numerical bound and realization inputs are not proved here.
\end{theorem}

\begin{proof}
\proofstep{1. The original host, coefficient and strict group.}
The original host contains $K,t,\zeta_7$ and the independent
finite factors of (FC.1). The two faithful Kummer powers and
six powers of $\xi_7$ are actual labels in that host; no new
finite character is adjoined. Their values lie in $\mu_6$,
contained in every prescribed original completion $O_7$ of
$\mathbb Z[\zeta_{3\cdot2^h}]$. Its original fraction field,
ramification, embedding and rank $b_7$ are retained. Literally,
\[
 (O_7\otimes_{\mathbb Z_7}\mathbb Q_7/\mathbb Z_7)
                         (\chi^j\xi_7^a)(-e)
       =(E_7/O_7)(\chi^j\xi_7^a c_7^{-e}).
 \tag{MRTW.13}
\]
Apply the full original-coefficient and actual-host equality in
G219, (MRT.16), (MRT.27) and Step7. Every extra host inertia
kernel is retained. Equality is essential because the new terms
have a specified nonzero size; a one-sided injection would not
establish that size. Full decomposition zero at both seven places
and inertia zero elsewhere are imposed on the original coefficient.
For matched $a,e$, uniqueness of Teichm\"uller lifts gives
$\xi_7^a=\omega_7^e$, so the action is
$\chi^j\langle c_7\rangle^{-e}$ on the same ring.
G219's finite-free original-$O_7$ comparison uses $O_7/7^nO_7$,
without replacing it by a power of the original uniformizer with
a different exponent. Consequently the full strict group has
order $7^{b_7s_e}$, and the normalized logarithm is $s_e\log7$.
The same $s_e$ serves both cubic labels by conjugation and every
host by the literal reverse-inflation equality. The groups are
finite at every positive weight by the retained original finiteness
input; the nested-residue conclusion (MRT.35)--(MRT.36) assumes no bound for
the largest such length at the cutoff.
The finite-etale Artin--Tate dictionary and complete applied
arithmetic proof remain those of G219. No new class computation
or independent primary-source proof is claimed here.

\proofstep{2. Every old intersection and the universal counts.}
Every new label has nontrivial Kummer factor and therefore avoids
all rational old sets $V_u,Y_u,C_5^{\rm all},C_7^{\rm all}$.
The mixed part of $Z_u$ occurs only at coefficient $269$.
The growing mixed families $M_8,M_7,C,A_{\rm ur}$ exclude seven.
At coefficient seven the old $Q_7$ has a nontrivial two-power
quadratic factor, unlike $B_u$. The old $Q_{7,7}$ has residual
order three. The part of $A_7$ with trivial two-power factor has
residual order two; its residual-order-one part requires a
nontrivial two-power quadratic factor. These exhaust $M_{C7}$
by (MQSW.4)--(MQSW.6), and all are disjoint from $B_u$.
The finite factor distinctions persist on $G_K$ by (FC.1),
independently of coincident primes or weight progressions.

Reduction on $\mu_6$ is injective at seven. For every positive
weight exactly one of the six powers of the actual faithful
$\xi_7$ has residual match, in each original embedding. The
independent two Kummer powers give two distinct triples at each
weight. This proves $|B_u|=2N$, $|P_m|=2$ and $|A_m|=2(N-1)$.
Add that disjoint count to (MQSW.8) to obtain (MRTW.10).
For $N\geq2$ each deletion is nonempty. For $N=1$ the unique
branch deletes no matched term and equals the old remainder.
These are counts for the full family in one theorem; neither
the former six-choice formulation nor smaller residue subfamilies
are staged as separate arithmetic reductions.

\proofstep{3. One maximal weight controls the entire discarded cost.}
All sums are finite by (CP.3)--(CP.5) and G192's Step3, retained
by G218. Separate the identical core from the matched terms.
Step1 then proves all three identities (MRTW.6).
Choose a weight $m_*$ attaining $s_*:=\max_{1\leq e\leq N}s_e$;
this choice is used only in the proof of the pointwise bound.
For $n=1$, there are $N-1$ other matched weights of length at
least one. For $2\leq n\leq s_*$, the unique possible class
modulo $7^{n-1}$ containing all weights of length at least $n$
contains $m_*$ itself. An interval of $N$ consecutive integers
has at most $1+\lfloor(N-1)/7^{n-1}\rfloor$ members in one
such class. Removing $m_*$ leaves at most
$\lfloor(N-1)/7^{n-1}\rfloor$. For $n>s_*$ none remain.
The finite layer-counting identity therefore gives
\[
 \begin{split}
 \sum_{\substack{1\leq e\leq N\\e\ne m_*}}s_e
 &=\sum_{n\geq1}
        \#\{1\leq e\leq N:e\ne m_*,\ s_e\geq n\}\\
 &\leq (N-1)+\sum_{k\geq1}
                      \left\lfloor\frac{N-1}{7^k}\right\rfloor
   =L_N\leq\frac76(N-1).
 \end{split}\tag{MRTW.14}
\]
There are two cubic labels of each length. Multiplying by
$2\log7$ proves the upper bound in (MRTW.8). The lower bound
follows from $s_e\geq1$ for every other weight. This also proves
strict properness of the maximum for $T\geq3$ and equality at
$T=2$. The upper bound does not contain $s_*$: no estimate of
the retained largest group has been inserted.
Since $\log7<3\log2$ and $\mathcal P_u(T)\geq T\log2$,
\[
 \frac73(N-1)\log7<7T\log2\leq7\mathcal P_u(T).
 \tag{MRTW.15}
\]
Thus the explicit maximum has a single known comparison cost,
although its number of branches grows with the actual cutoff.

\proofstep{4. One old comparison and the unchanged exponent.}
Apply the triangle inequality to the two separate comparisons
(MQSW.9) and the exact nonnegative difference (MRTW.8).
This proves (MRTW.11). Their full-local errors are not added
together. The majorant (FL.4) was already summed over all
original triples, and (EU.2) and (LT.2) retain every positive
derived index and original degree-one torsion term.
No index is declared zero or split into separately positive pieces.
The critical partners of the old rational noncritical sum remain
inside the same $C_2\mathcal P_u$ aggregate. They are neither
removed nor charged again by the matched mixed family.
There is no factor $N$ in the comparison cost: (MRTW.6)--(MRTW.8)
compare the pointwise maximum directly to one old remainder,
without summing $N$ separate error estimates.
The degrees $D_u=144\cdot2^h$ and $d_u=96\cdot2^h$ keep their
separate roles in (CAW.5). Replacing $C_2$ by $C_2+7$ in
(CAW.7), using (CP.5), proves the unchanged envelope (MRTW.12).

\proofstep{5. The same window with the required slack adjustment.}
For any real threshold $B$ the exact equivalence is
\[
 \mathcal R^{\rm old}(u)\leq B
 \quad\Longleftrightarrow\quad
 \mathcal R_{399}^{\rm matched,max}(u)\leq B-\Gamma_u.
 \tag{MRTW.16}
\]
A uniform sufficient condition uses $B-2L_N\log7$; the bound
$B-7\mathcal P_u$ is also sufficient. The old bound also implies the new maximum
bound at threshold $B$, by nonnegativity. No unchanged numerical
slack at a fixed finite threshold is asserted.

Fix the original finite datum, page and $0<\varepsilon<1/256$.
Assume the exact (TC.3) quantifiers with the new remainder:
there are $J\geq1$, $X_0$ and $s\in(0,1)$, independent of
$X,q,d,u,\ell,\eta,e$, such that for every original prescribed
pair at $X\geq X_0$ one window satisfies
\[
 X^{1/J}\leq W,\qquad W^2\leq X,\qquad
 \mathcal R_{399}^{\rm matched,max}(u)
             \leq(1-s)\exp(u^\varepsilon/2)
                       \quad(W\leq u\leq W^2).
 \tag{MRTW.17}
\]
Choose a uniform $U_0$ so that the full error in the first line
of (MRTW.11), with its enlarged constant, is at most
$(s/2)\exp(u^\varepsilon/2)$ for all permitted $u\geq U_0$.
Increase the source threshold to $\max(X_0,U_0^J)$ and keep
that identical $W$. The full original numerical index premise
holds on the same interval with slack $s/2$.
Starting instead from the preceding $I_{399}^{\rm all7,2}$
premise, the second comparison gives the new remainder premise
by the same threshold and half-slack argument.
Every branch $1\leq m\leq N(u)$ is tested on this one $W$.
A maximizing weight may vary with $u$; it need not be chosen
uniformly across the interval. The maximum itself is defined
before any window is selected, and its cost bound is uniform.
No separately existential weight-dependent windows are intersected.
Uniform little-$o$ versions and fixed finite simultaneous families
use a maximum of their uniform error constants, retaining every
original rank $R.t$, lattice, page, pair and realization map.
\end{proof}

\paragraph{Scope and chronology.}
This is a controlled cost of nonzero matched cubic groups, not
strict vanishing or a formula for individual positive indices.
One largest matched pair, every other unsupported mixed term
and the entire retained noncritical rational sum remain in the
explicit maximum. Their numerical upper bound, the mixed
realization $\mathcal R_{\rm cub}$ and H04.4 remain open.
The complete higher-lifting theorem and this one composition
supersede the earlier six-choice WIP without a separate final
stage or reserved later movement credit. Every completion and
review time is recorded only when actually reached.
